\chapter{S1-B 빠른 Sorting \& Iterator 참고}

\section{Iterator category 빠른 표}
\begin{longtable}{p{3.3cm}p{10.2cm}}
\toprule
\textbf{Category} & \textbf{대표 능력} \\
\midrule
Input & 읽기 \code{*it}, 전진 \code{++it}; single-pass 가능 \\
Output & 쓰기 \code{*it=value}, 전진 \code{++it} \\
Forward & multi-pass forward traversal \\
Bidirectional & forward + \code{--it} \\
Random-access & \code{+n}, \code{-n}, \code{+=}, \code{-=}, \code{it[n]}, iterator subtraction/order comparison \\
Contiguous (C++20) & random-access + contiguous storage \\
\bottomrule
\end{longtable}

\section{Container별 iterator}
\begin{longtable}{p{5cm}p{8.5cm}}
\toprule
\textbf{Container/Source} & \textbf{Iterator 성질} \\
\midrule
\code{vector}, \code{array}, \code{string} & random-access; contiguous storage \\
\code{deque} & random-access, non-contiguous \\
\code{list} & bidirectional \\
\code{forward\_list} & forward \\
\code{set}, \code{map}, \code{multiset}, \code{multimap} & bidirectional \\
unordered associative containers & forward \\
\code{istream\_iterator} & input \\
\code{ostream\_iterator} & output \\
\bottomrule
\end{longtable}

\section{S1-B 핵심 연산}
\begin{lstlisting}
auto it = v.begin();

*it;          // dereference
++it;         // forward
--it;         // bidirectional or stronger

it + n;       // random-access
it - n;
it += n;
it -= n;
it[n];
b - a;        // iterator distance, random-access
\end{lstlisting}

\section{\code{iterator} vs \code{const\_iterator}}
\begin{lstlisting}
vector<int> v = {1, 2, 3};

auto it = v.begin();
*it = 10; // allowed

auto cit = v.cbegin();
// *cit = 10; // error
\end{lstlisting}

\section{Generic utilities}
\begin{lstlisting}
auto a = next(it, 3);
auto b = prev(it, 2);
advance(it, 5);
auto d = distance(first, last);
\end{lstlisting}

Random-access iterator에서는 여러 칸 이동과 distance가 $O(1)$인 반면, forward/bidirectional iterator에서는 $O(n)$일 수 있다.

\section{\code{vector}와 \code{set}의 중요한 차이}
\begin{lstlisting}
// vector: random-access
auto count = right - left;

// set: bidirectional
// right - left; // error
auto count2 = distance(left, right); // potentially O(N)
\end{lstlisting}

\section{Ordered container의 bound}
\begin{lstlisting}
set<int> s;

auto it1 = s.lower_bound(x); // O(log N)

// Generic lower_bound on set iterators can incur
// linear iterator movement.
auto it2 = lower_bound(s.begin(), s.end(), x);
\end{lstlisting}

\section{Half-open range}
\[
[\code{begin()},\code{end()})
\]
\code{end()}는 마지막 원소 다음 위치이므로 dereference하지 않는다.

\section{\code{map} iterator}
\begin{lstlisting}
map<int, string> mp;
auto it = mp.begin();

// it->first = 7;    // error
it->second = "new";  // allowed
\end{lstlisting}

\code{map<K,V>::value\_type}은 \code{pair<const K,V>}다.

\section{Iterator invalidation 빠른 점검}
\begin{itemize}
\item \code{vector}: reallocation 시 전체 iterator 무효화; insert/erase 위치 이후도 주의
\item \code{list}/\code{map}/\code{set}: insert는 대체로 기존 iterator 보존; erase된 원소 iterator는 무효
\item unordered container: rehash 시 iterator 무효화 가능
\end{itemize}
